\subsection{非紧区间上积分的定义}

设 I 是紧区间 \([a, b]\)，含于扩充实直线 \(\overline{R}\)（a 与 b 可以是无穷）；设 f 是定义在 {]}a, b{[} 上、取值于 R 上的完备赋范空间 E 中的函数。把 II, p.~51 的定义 1 加以推广，我们说：定义在 \([a, b]\) 上且取值于 E 中的函数 g 是 f 的原函数，如果它在 \([a, b]\) 上连续（特别是在端点 a 与 b 处连续），并且在这区间的一个可数子集的余集的各点处具有等于 \(\mathbf{f}(x)\) 的导数（余集相对于 {]}a, b{[} 取）。

我们将限于如下情形：存在一个有限的严格递增序列 \((c_i)_{0\leqslant i\leqslant n}\)，由 \(\mathrm{I} = [a,b]\) 的点组成，使得 \(c_0 = a\)、\(c_n = b\)，并且使得 \(\mathbf{f}\) 在每个开区间 \(]c_i\) , \(c_{i + 1}[\) 上是规则的，尽管在包含至少一个点 \(c_i\)（作为 \(\mathrm{I}\) 的内点）的每个开区间上不一定都是规则的；这样的函数称为 \(]a\) , \(b[\) 上的分段规则函数。我们注意，\(]a\) , \(b[\) 上的规则函数是分段规则的（在前面的定义中取 \(n = 1\)）。

若 \(\mathbf{f}\)（按上面明确的意义）有一个原函数 \(\mathbf{g}\)，且若 \(c\) 是区间 \([c_i, c_{i+1}[(0 \leqslant i \leqslant n-1)\) 中的一点，则据假设，对此区间中的每个 \(x\)，有 \(\mathbf{g}(x) - \mathbf{g}(c) = \int_c^x \mathbf{f}(t) dt\)；由于 \(\mathbf{g}\) 据假设在 I 上连续，可见 \(\int_c^x \mathbf{f}(t) dt\) 当 \(x\) 从右边趋于 \(c_i\) 时以及当 \(x\) 从左边趋于 \(c_{i+1}\) 时在 E 中趋于一个极限。反之，假设这些条件对一切 \(t\) 都满足，并设 \(\mathbf{g}_i\) 是 \(\mathbf{f}\) 在区间 \([c_i, c_{i+1}[(0 \leqslant i \leqslant n-1)\) 上的一个原函数；我们立即注意到：函数 \(\mathbf{g}\) 在诸点 \(c_i\) 所成之集相对于 I 的余集上按如下条件定义：它等于 \(\mathbf{g}_i(x) + \sum_{k=1}^i (\mathbf{g}_{k-1}(c_k-) - \mathbf{g}_k(c_k+))\)（在 \([c_i, c_{i+1}[\) 上，对 \(0 \leqslant i \leqslant n-1\)）；它在 I 的每个异于诸 \(c_i\) 的点处连续，并且在每个这样的点处有极限；因此它可以连续地延拓到每个这样的点，而延拓所得的函数显然是 \(\mathbf{f}\) 在 I 上的一个原函数。此外显然，\(\mathbf{f}\) 的任何其他原函数都具有形式 \(\mathbf{g} + \mathbf{a}\)（\(\mathbf{a}\) 是 E 的一个元素）。

定义 1. 我们说向量函数 \(\mathbf{f}\)（在区间 \(]a\) , \(b[\)（属于 \(\overline{\mathbf{R}}\)）上分段规则的）在这个区间上有一个积分，如果 \(\mathbf{f}\) 在 \([a, b]\) 上有一个原函数；若 \(\mathbf{g}\) 是 \(\mathbf{f}\) 在 \([a, b]\) 上的任一原函数，且 \(x_0\) 与 \(x\) 是 \([a, b]\) 上的任意两点，则称元素 \(\mathbf{g}(x) - \mathbf{g}(x_0)\) 为 \(\mathbf{f}\) 从 \(x_0\) 到 \(x\) 的积分，并记作 \(\int_{x_0}^x \mathbf{f}(t) dt\)。

当区间 \([x_0, x]\) 不含任何点 \(c_i\) 时，这个概念显然与此前所定义的概念一致。

定义 1 之前的这些注记表明，为使 \(\mathbf{f}\) 在 \(]a\) , \(b[\) 上有一个积分，必须且只需它在每个区间 \([c_i, c_{i+1}]\) 上的限制在该区间上有一个积分。换言之，可以化归为如下的情形：\(\mathbf{f}\) 在非紧区间 \(\mathrm{I} \subset \mathbf{R}\) 上是规则的，其端点为 \(a, b\)（\(a < b\)），并且：\(1^\circ\) 数 \(a, b\) 中（至少）有一个是无穷的；\(2^\circ\) 或者 \(\mathbf{f}\) 在包含点 \(a, b\) 中至少一个的每个紧区间上不是规则的（这两个假设并不互斥）。为使 \(\mathbf{f}\) 在 \(\mathrm{I}\) 上有一个积分，必须且只需积分 \(\int_x^y \mathbf{f}(t) dt\) 当点 \((x, y)\) 趋于 \((a, b) \in \overline{\mathbf{R}}^2\) 而保持属于 \(\mathrm{I} \times \mathrm{I}\) 时趋于一个极限：而这个极限不是别的，正是依定义 1 的 \(\int_a^b \mathbf{f}(t) dt\)。滥用语言的说法，人们不说 \(\mathbf{f}\) 在 \(\mathrm{I}\) 上有一个积分，而说积分 \(\int_a^b \mathbf{f}(t) dt\) 是收敛的（或收敛）。

例。1) 积分 \(\int_{1}^{+\infty} dt / t^{2}\) 收敛且等于 1，因为

\[
\int_ {1} ^ {x} \frac {d t}{t ^ {2}} = 1 - \frac {1}{x}.
\]

\begin{enumerate}
\def\labelenumi{\arabic{enumi})}
\setcounter{enumi}{1}
\tightlist
\item
  积分 \(\int_0^1 dt / \sqrt{t}\) 收敛且等于 2，因为
\end{enumerate}

\[
\int_ {x} ^ {1} \frac {d t}{\sqrt {t}} = 2 (1 - \sqrt {x}) \quad \text { for } \quad x > 0.
\]

\begin{enumerate}
\def\labelenumi{\arabic{enumi})}
\setcounter{enumi}{2}
\tightlist
\item
  设 \((\mathbf{u}_n)_{n\geqslant 1}\) 是 E 的点的一个无穷序列，\(\mathbf{f}\) 是由下列条件定义在区间 {[}1, \(+\infty[\) 上的阶梯函数：\(\mathbf{f}(x) = \mathbf{u}_n\)（当 \(n\leqslant x < n + 1\)）。那么，为使积分 \(\int_1^{+\infty}\mathbf{f}(t)dt\) 收敛，必须且只需通项为 \(\mathbf{u}_n\) 的级数在 E 中收敛；事实上，有
\end{enumerate}

\[
\int_ {1} ^ {n} \mathbf {f} (t) d t = \sum_ {p = 1} ^ {n - 1} \mathbf {u} _ {p},
\]

故条件是必要的；反之，若通项为 \(\mathbf{u}_n\) 的级数在 E 中收敛，则 \(\lim_{n\to \infty}\mathbf{u}_n = 0\)；而今若 \(n\leqslant x\leqslant n + 1\)，就有 \(\int_1^n\mathbf{f}(t)dt = \sum_{p = 1}^{n - 1}\mathbf{u}_p +\) \(\mathbf{u}_n(x - n)\)，故这个积分具有极限 \(\sum_{n = 1}^{\infty}\mathbf{u}_n\)（当 \(x\) 趋于 \(+\infty\) 时）。

立即可以看出，若分段规则函数 \(\mathbf{f}\) 在 I 上有一个积分，则 II, p.~59 的公式 (4) 至 (9) 仍然有效。类似地，II, p.~59 的公式 (10) 按如下方式推广：假设 \(\mathbf{f}\) 与 \(\mathbf{g}\) 是规则函数 \(\mathbf{f}'\) 与 \(\mathbf{g}'\) 在 \(]a\) , \(b[\) 上的原函数，并用 \([\mathbf{f}.\mathbf{g}]|_a^b\) 表示 \([\mathbf{f}.\mathbf{g}]|_x^y\) 当 \((x, y)\) 趋于 \((a, b)\)（且 \(a < x \leqslant y < b\)）时的极限（如果它存在）；那么，若（三个）表达式 \([\mathbf{f}.\mathbf{g}]|_a^b\)、\(\int_a^b [\mathbf{f}(t).\mathbf{g}'(t)] dt\) 与 \(\int_a^b [\mathbf{f}'(t).\mathbf{g}(t)] dt\) 中有两个有意义，则第三个也有意义，并且 II, p.~59 的公式 (10) 成立。

最后，设 \(f\) 是定义且连续在 \(I = ]a\) , \(b[\) 上的实函数，并且是某个规则函数 \(f'\) 在 \(]a\) , \(b[\) 上的原函数；另一方面，设 \(\mathbf{g}\) 是开区间 \(J\)（包含 \(f(I)\)）上的连续向量函数；若函数 \(\mathbf{g}(f(x))f'(x)\) 在 \(I\) 上有一个积分，且若 \(f\) 在点 \(a\) 与 \(b\) 处趋于一个极限（有限或无穷），则 \(\mathbf{g}\) 有一个从 \(f(a+)\) 到 \(f(b-)\) 的积分，并且有公式

\[
\int_ {a} ^ {b} \mathbf {g} (f (t)) f ^ {\prime} (t) d t = \int_ {f (a +)} ^ {f (b -)} \mathbf {g} (u) d u.\tag{1}
\]

事实上，若 \((x, y)\) 趋于 \((a, b)\)，则据假设 \((f(x), f(y))\) 趋于 \((f(a+), f(b-))\)；只需在 x 与 y 之间应用 II, p.~60 的公式 (12) 并取极限，即得 (1)。

给定非紧区间 \(I \subset R\) 上的一个规则函数 f，其端点为 a 与 b（a \textless{} b），f 在 I 上有积分这一条件可以表述如下。紧区间 \(J \subset I\) 构成一个有向集 \(\mathfrak{K}(I)\)（关于关系 \(\subset^{1}\)），因为若 \([\alpha, \beta]\) 与 \([\gamma, \delta]\) 是含于 I 的两个紧区间，且令 \(\lambda = \min(\alpha, \gamma)\)、\(\mu = \max(\beta, \delta)\)，则区间 \([\lambda, \mu]\) 含于 I 且包含所考虑的两个区间。对每个含于 I 的紧区间 \(J = [\alpha, \beta]\)，令

\[
\int_ {\mathrm{J}} \mathbf {f} (t) d t = \int_ {\alpha} ^ {\beta} \mathbf {f} (t) d t;
\]

为使 f 在 I 上有一个积分，必须且只需映射 \(\mathbf{J} \mapsto \int_{\mathrm{J}} \mathbf{f}(t) dt\) 关于有向集 \(\mathcal{K}(\mathrm{I})\) 有一个极限；这个极限就是积分 \(\int_{a}^{b} \mathbf{f}(t) dt\)，我们仍把它记作 \(\int_{\mathrm{I}} \mathbf{f}(t) dt\)。

命题 1（积分的柯西准则）。设 \(\mathbf{f}\) 是区间 \(\mathrm{I} \subset \mathbf{R}\) 上的规则函数，其端点为 \(a\) 与 \(b\)（\(a < b\)）。为使积分 \(\int_{a}^{b} \mathbf{f}(t) dt\) 存在，必须且只需：对每个 \(\varepsilon > 0\)，存在一个含于 I 的紧区间 \(\mathrm{J}_0 = [\alpha, \beta]\)，使得对任何紧区间 \(\mathrm{K} = [x, y]\)（含于 I 且与 \(\mathrm{J}_0\) 无公共内点），都有 \(\left\| \int_{\mathrm{K}} \mathbf{f}(t) dt \right\| \leqslant \varepsilon\)。

事实上，由于 E 是完备的，柯西准则表明：为使积分 \(\int_{\mathrm{I}} \mathbf{f}(t) dt\) 收敛，必须且只需对任何 \(\varepsilon > 0\) 存在一个紧区间 \(\mathrm{J}_0 = [\alpha, \beta]\)，使得对每个满足 \(\mathrm{J}_0 \subset \mathrm{J} \subset I\) 的紧区间 J，有 \(\left\| \int_{\mathrm{J}} \mathbf{f}(t) dt - \int_{\mathrm{J}_0} \mathbf{f}(t) dt \right\| \leqslant \varepsilon\)。该命题将由下面的引理推出：

引理。设 \(J_{0} = [\alpha, \beta]\) 是含于 I 的紧区间。为使 \(\left\|\int_{J} \mathbf{f}(t) dt - \int_{J'} \mathbf{f}(t) dt\right\| \leqslant \varepsilon\) 对每一对含于 I 且包含 \(J'\) 的紧区间 J、\(J'\) 都成立，必要的是 \(\left\|\int_{K} \mathbf{f}(t) dt\right\| \leqslant \varepsilon\)，而且 \(\left\|\int_{K} \mathbf{f}(t) dt\right\| \leqslant \varepsilon/2\) 就是充分的（对每个含于 I 且与 \(J_{0}\) 无公共内点的紧区间 K）。

事实上，若对 \(J_{0} \subset J \subset I\) 与 \(J_{0} \subset J' \subset I\) 有

\[
\left\| \int_ {J} \mathbf {f} (t) d t - \int_ {J ^ {\prime}} \mathbf {f} (t) d t \right\| \leqslant \varepsilon
\]

则特别地可以看出，对 \(x \leqslant y \leqslant \alpha\)，或对 \(\beta \leqslant x \leqslant y\)（x 与 y 属于 I），有 \(\left\|\int_{x}^{y} \mathbf{f}(t) dt\right\| \leqslant \varepsilon\)。反之，若 \(\left\|\int_{K} \mathbf{f}(t) dt\right\| \leqslant \varepsilon/2\) 对每个紧区间 \(K \subset I\)（满足 \(K \cap J_{0} = \emptyset\)）都成立，且若 \(J = [x, y]\)、\(J' = [z, t]\) 是含于 I 且包含 \(J_{0}\) 的两个紧区间，则有

\[
\left\| \int_ {J} \mathbf {f} (t) d t - \int_ {J ^ {\prime}} \mathbf {f} (t) d t \right\| = \left\| \int_ {x} ^ {z} \mathbf {f} (t) d t + \int_ {t} ^ {y} \mathbf {f} (t) d t \right\| \leqslant \varepsilon ,
\]

因为

\[
x \leqslant \alpha \leqslant \beta \leqslant y \quad \text { and } \quad z \leqslant \alpha \leqslant \beta \leqslant t.
\]

例。若区间 I 有界，且 \(\mathbf{f}\) 在 I 上有界，则积分 \(\int_{\mathrm{I}}\mathbf{f}(t)dt\) 总存在，因为由中值定理，对 \(y\leqslant \alpha \leqslant \beta \leqslant z\) 有

\[
\left\| \int_ {y} ^ {\alpha} \mathbf {f} (t) d t \right\| \leqslant (\alpha - a) \sup _ {x \in I} \| \mathbf {f} (x) \|, \quad \left\| \int_ {\beta} ^ {z} \mathbf {f} (t) d t \right\| \leqslant (b - \beta) \sup _ {x \in I} \| \mathbf {f} (x) \|
\]

只需取 \(\alpha - a\) 与 \(b - \beta\) 足够小，柯西准则便得到满足。可以注意，在这种情形下，f 在 I 上的原函数在 I 的左端点（相应地，右端点）处不一定有右（相应地，左）导数（当该数为有限时），这与 I 紧且 f 在 I 上规则的情形不同（参见 II, p.~33, 习题 1）。

